\section{Boolean model integer gap}\label{sec:boolean_model_appendix}

The integer gap is defined as
\[
    \frac{\text{Optimal value of the IP} - \text{Optimal value of the LP relaxation}}{\text{Optimal value of the LP relaxation}}.
\]
The optimal value of the LP relaxation is a lower bound for the optimal value of an IP\@.
When pricing completed, which is true for green cells that do not say ``Limit'',
the LP relaxation value observed is the true optimal LP relaxation value.
Hence, no integer solution with an objective value better than it exists.
\autoref{tab:PDL1ex1_solve_time} shows that pricing finished for all instances, and this argument holds.
\autoref{tab:PDL1ex1neg_solve_time} shows that pricing only finished in some cases.
For the cases in which pricing completed, the gap is small.
Hence, for these cases, we conclude that there do not exist meaningfully better rulesets than the ones we found.
It also implies branch-and-price would not be meaningfully better than price-then-branch in this case.
For the cases where pricing did not converge, this conclusion cannot be made.

\begin{table}
    \begin{tabular}{lrrrrrr}
\toprule
\textbf{$n_{\text{quantiles}}$} & \textbf{1} & \textbf{3} & \textbf{5} & \textbf{7} & \textbf{9} \\
\textbf{$Q_y$} &  &  &  &  &  \\
\midrule
\textbf{0.5} & {\cellcolor[HTML]{C6E0B4}} \num{0.000} & {\cellcolor[HTML]{C6E0B4}} \num{0.000} & {\cellcolor[HTML]{C6E0B4}} \num{0.004} & {\cellcolor[HTML]{C6E0B4}} \num{0.010} & {\cellcolor[HTML]{C6E0B4}} \num{0.003} \\
\textbf{0.6} & {\cellcolor[HTML]{C6E0B4}} \num{0.000} & {\cellcolor[HTML]{C6E0B4}} \num{0.012} & {\cellcolor[HTML]{C6E0B4}} \num{0.000} & {\cellcolor[HTML]{C6E0B4}} \num{0.001} & {\cellcolor[HTML]{C6E0B4}} \num{0.000} \\
\textbf{0.7} & {\cellcolor[HTML]{C6E0B4}} \num{0.000} & {\cellcolor[HTML]{C6E0B4}} \num{0.000} & {\cellcolor[HTML]{C6E0B4}} \num{0.000} & {\cellcolor[HTML]{C6E0B4}} \num{0.002} & {\cellcolor[HTML]{C6E0B4}} \num{0.000} \\
\textbf{0.8} & {\cellcolor[HTML]{C6E0B4}} \num{0.000} & {\cellcolor[HTML]{C6E0B4}} \num{0.005} & {\cellcolor[HTML]{C6E0B4}} \num{0.000} & {\cellcolor[HTML]{C6E0B4}} \num{0.002} & {\cellcolor[HTML]{C6E0B4}} \num{0.000} \\
\textbf{0.9} & {\cellcolor[HTML]{C6E0B4}} \num{0.000} & {\cellcolor[HTML]{C6E0B4}} \num{0.000} & {\cellcolor[HTML]{C6E0B4}} \num{0.000} & {\cellcolor[HTML]{C6E0B4}} \num{0.000} & {\cellcolor[HTML]{C6E0B4}} \num{0.000} \\
\bottomrule
\end{tabular}

    \caption{Integer gap for the boolean model on the full Broad dataset for various $\nq$ and $Q_y$,
        only lower-bounds}
    \label{tab:PDL1ex1_solve_time}
\end{table}

\begin{table}
    \begin{tabular}{lrrrrrr}
\toprule
\textbf{$n_{\text{quantiles}}$} & \textbf{1} & \textbf{3} & \textbf{5} & \textbf{7} & \textbf{9} \\
\textbf{$Q_y$} &  &  &  &  &  \\
\midrule
\textbf{0.5} & {\cellcolor[HTML]{C6E0B4}} \num{0.028} & {\cellcolor[HTML]{F4CCCC}} Limit & {\cellcolor[HTML]{F4CCCC}} Limit & {\cellcolor[HTML]{F4CCCC}} Limit & {\cellcolor[HTML]{F4CCCC}} Limit \\
\textbf{0.6} & {\cellcolor[HTML]{C6E0B4}} \num{0.015} & {\cellcolor[HTML]{F4CCCC}} Limit & {\cellcolor[HTML]{F4CCCC}} Limit & {\cellcolor[HTML]{F4CCCC}} Limit & {\cellcolor[HTML]{F4CCCC}} Limit \\
\textbf{0.7} & {\cellcolor[HTML]{C6E0B4}} \num{0.003} & {\cellcolor[HTML]{C6E0B4}} \num{0.009} & {\cellcolor[HTML]{F4CCCC}} Limit & {\cellcolor[HTML]{F4CCCC}} Limit & {\cellcolor[HTML]{F4CCCC}} Limit \\
\textbf{0.8} & {\cellcolor[HTML]{C6E0B4}} \num{0.000} & {\cellcolor[HTML]{C6E0B4}} \num{0.009} & {\cellcolor[HTML]{F4CCCC}} Limit & {\cellcolor[HTML]{F4CCCC}} Limit & {\cellcolor[HTML]{F4CCCC}} Limit \\
\textbf{0.9} & {\cellcolor[HTML]{C6E0B4}} \num{0.000} & {\cellcolor[HTML]{C6E0B4}} \num{0.000} & {\cellcolor[HTML]{C6E0B4}} \num{0.011} & {\cellcolor[HTML]{C6E0B4}} \num{0.000} & {\cellcolor[HTML]{F4CCCC}} Limit \\
\bottomrule
\end{tabular}

    \caption{Integer gap for the boolean model on the full Broad dataset for various $\nq$ and $Q_y$,
        both lower-bounds and upper-bounds}
    \label{tab:PDL1ex1neg_solve_time}
\end{table}


%\begin{table}[h]
%    \centering
%    \begingroup
%    \renewcommand{\cline}[1]{}
%    \rowcolors{2}{gray!10}{white}
%    \arrayrulecolor{black}
%    \begin{tabular}{lllrrrl}
\toprule
 &  &  & \textbf{Acc.} & \textbf{Prec.} & \textbf{Rec.} & \textbf{Ruleset} \\
$\mathbf{Q_y}$ & $\mathbf{C}$ & $\mathbf{D}$ &  &  &  &  \\
\midrule
\multirow[t]{3}{*}{0.5} & \multirow[t]{3}{*}{1} & 1 & 0.71 & 0.68 & 0.81 & \makecell{\{$\text{FOSL1} \geq 30.61 \text{ Q}(0.40)$\}} \\
 &  & 2 & 0.73 & 0.75 & 0.70 & \makecell{\{$\text{FOSL1} \geq 30.61 \text{ Q}(0.40)$ $\wedge$ \\$\text{STAT4} \geq 0.17 \text{ Q}(0.31)$\}} \\
 &  & 3 & 0.74 & 0.76 & 0.70 & \makecell{\{$\text{FOSL1} \geq 30.61 \text{ Q}(0.40)$ $\wedge$ \\$\text{STAT4} \geq 0.17 \text{ Q}(0.31)$ $\wedge$ \\$\text{NFATC1} \leq 9.69 \text{ Q}(0.90)$\}} \\
\cline{1-7} \cline{2-7}
\multirow[t]{3}{*}{0.6} & \multirow[t]{3}{*}{1} & 1 & 0.74 & 0.68 & 0.68 & \makecell{\{$\text{FOSL1} \geq 109.35 \text{ Q}(0.60)$\}} \\
 &  & 2 & 0.76 & 0.75 & 0.60 & \makecell{\{$\text{FOSL1} \geq 109.35 \text{ Q}(0.60)$ $\wedge$ \\$\text{NFATC1} \leq 3.92 \text{ Q}(0.70)$\}} \\
 &  & 3 & 0.76 & 0.75 & 0.60 & \makecell{\{$\text{AR} \geq 0.00 \text{ Q}(0.28)$ $\wedge$ \\$\text{FOSL1} \geq 109.35 \text{ Q}(0.60)$ $\wedge$ \\$\text{NFATC1} \leq 3.92 \text{ Q}(0.70)$\}} \\
\cline{1-7} \cline{2-7}
\bottomrule
\end{tabular}

%    \endgroup
%    \caption{}
%    \label{tab:tpm_interpretable_results_C_1}
%\end{table}
%
%\begin{table}[h]
%    \centering
%    \begingroup
%    \renewcommand{\cline}[1]{}
%    \rowcolors{2}{gray!10}{white}
%    \arrayrulecolor{black}
%    \begin{tabular}{lllrrrl}
\toprule
 &  &  & \textbf{Acc.} & \textbf{Prec.} & \textbf{Rec.} & \textbf{Ruleset} \\
$\mathbf{Q_y}$ & $\mathbf{C}$ & $\mathbf{D}$ &  &  &  &  \\
\midrule
\multirow[t]{3}{*}{0.5} & \multirow[t]{3}{*}{2} & 1 & 0.71 & 0.68 & 0.81 & \makecell{\{$\text{FOSL1} \geq 30.61 \text{ Q}(0.40)$\}} \\
 &  & 2 & 0.76 & 0.79 & 0.71 & \makecell{\{$\text{FOSL1} \geq 109.35 \text{ Q}(0.60)$ $\wedge$ \\$\text{NFATC1} \leq 3.92 \text{ Q}(0.70)$\}\\ $\vee$ \\\{$\text{IRF1} \geq 12.67 \text{ Q}(0.50)$ $\wedge$ \\$\text{MYC} \leq 122.33 \text{ Q}(0.50)$\}} \\
 &  & 3 & 0.76 & 0.82 & 0.67 & \makecell{\{$\text{FOSL1} \geq 109.35 \text{ Q}(0.60)$ $\wedge$ \\$\text{MYC} \geq 54.30 \text{ Q}(0.20)$ $\wedge$ \\$\text{NFATC1} \leq 3.92 \text{ Q}(0.70)$\}\\ $\vee$ \\\{$\text{IRF1} \geq 12.67 \text{ Q}(0.50)$ $\wedge$ \\$\text{STAT4} \geq 0.09 \text{ Q}(0.20)$ $\wedge$ \\$\text{MYC} \leq 122.33 \text{ Q}(0.50)$\}} \\
\cline{1-7} \cline{2-7}
\multirow[t]{3}{*}{0.6} & \multirow[t]{3}{*}{2} & 1 & 0.74 & 0.68 & 0.68 & \makecell{\{$\text{FOSL1} \geq 109.35 \text{ Q}(0.60)$\}} \\
 &  & 2 & 0.78 & 0.74 & 0.67 & \makecell{\{$\text{FOSL1} \geq 109.35 \text{ Q}(0.60)$ $\wedge$ \\$\text{NFATC1} \leq 3.92 \text{ Q}(0.70)$\}\\ $\vee$ \\\{$\text{STAT2} \geq 59.29 \text{ Q}(0.80)$ $\wedge$ \\$\text{STAT3} \geq 79.44 \text{ Q}(0.80)$\}} \\
 &  & 3 & 0.78 & 0.74 & 0.69 & \makecell{\{$\text{FOSL1} \geq 109.35 \text{ Q}(0.60)$ $\wedge$ \\$\text{NFATC1} \leq 3.92 \text{ Q}(0.70)$ $\wedge$ \\$\text{STAT2} \leq 48.42 \text{ Q}(0.70)$\}\\ $\vee$ \\\{$\text{IRF1} \geq 12.67 \text{ Q}(0.50)$ $\wedge$ \\$\text{STAT2} \geq 35.33 \text{ Q}(0.50)$ $\wedge$ \\$\text{MYC} \leq 122.33 \text{ Q}(0.50)$\}} \\
\cline{1-7} \cline{2-7}
\bottomrule
\end{tabular}

%    \endgroup
%    \caption{}
%    \label{tab:tpm_interpretable_results_C_2}
%\end{table}


% \begin{table}[h]
%  \centering
%  \begin{minipage}{0.1\linewidth}
%   \centering
%   \begin{tabular}{| c|c|}\hline
\textbf{Depth} & \multicolumn{1}{|>{\columncolor[HTML]{B8E6B8}}c|}{Baseline: } \\
 & \multicolumn{1}{|>{\columncolor[HTML]{B8E6B8}}c|}{0.69} \\ \hline
1 & \multicolumn{1}{|>{\columncolor[HTML]{C4DCB8}}c|}{-STAT2:} \\
 & \multicolumn{1}{|>{\columncolor[HTML]{C4DCB8}}c|}{0.69} \\ \hline
2 & \multicolumn{1}{|>{\columncolor[HTML]{EABFB8}}c|}{-RELA:} \\
 & \multicolumn{1}{|>{\columncolor[HTML]{EABFB8}}c|}{0.66} \\ \hline
3 & \multicolumn{1}{|>{\columncolor[HTML]{F5B8B8}}c|}{-FOSL1:} \\
 & \multicolumn{1}{|>{\columncolor[HTML]{F5B8B8}}c|}{0.66} \\ \hline
\end{tabular}
%  \end{minipage}
%  \hfill
%  \begin{minipage}{0.85\linewidth}
%   \centering
%   \begin{tabular}{||c|c|c|c|c|c|c|c|}\hline
\multicolumn{8}{|>{\columncolor[HTML]{B8E6B8}}c|}{Baseline: } \\
\multicolumn{8}{|>{\columncolor[HTML]{B8E6B8}}c|}{0.71} \\ \hline
\multicolumn{4}{|>{\columncolor[HTML]{BAE4B8}}c|}{-FOS:} & \multicolumn{4}{|>{\columncolor[HTML]{CDD5B8}}c|}{-STAT2:} \\
\multicolumn{4}{|>{\columncolor[HTML]{BAE4B8}}c|}{0.71} & \multicolumn{4}{|>{\columncolor[HTML]{CDD5B8}}c|}{0.69} \\ \hline
\multicolumn{2}{|>{\columncolor[HTML]{C0DFB8}}c|}{-RELA:} & \multicolumn{2}{|>{\columncolor[HTML]{CDD5B8}}c|}{-STAT2:} & \multicolumn{2}{|>{\columncolor[HTML]{E5C3B8}}c|}{-RELA:} & \multicolumn{2}{|>{\columncolor[HTML]{CFD3B8}}c|}{-SOX2:} \\
\multicolumn{2}{|>{\columncolor[HTML]{C0DFB8}}c|}{0.70} & \multicolumn{2}{|>{\columncolor[HTML]{CDD5B8}}c|}{0.69} & \multicolumn{2}{|>{\columncolor[HTML]{E5C3B8}}c|}{0.68} & \multicolumn{2}{|>{\columncolor[HTML]{CFD3B8}}c|}{0.69} \\ \hline
\multicolumn{1}{|>{\columncolor[HTML]{C4DCB8}}c|}{-NFKB1:} & \multicolumn{1}{|>{\columncolor[HTML]{E5C3B8}}c|}{-STAT2:} & \multicolumn{1}{|>{\columncolor[HTML]{E5C3B8}}c|}{-RELA:} & \multicolumn{1}{|>{\columncolor[HTML]{D2D2B8}}c|}{-SOX2:} & \multicolumn{1}{|>{\columncolor[HTML]{F5B8B8}}c|}{-FOSL1:} & \multicolumn{1}{|>{\columncolor[HTML]{EAC0B8}}c|}{-NFKB1:} & \multicolumn{1}{|>{\columncolor[HTML]{D2D2B8}}c|}{-FOS:} & \multicolumn{1}{|>{\columncolor[HTML]{E5C3B8}}c|}{-RELA:} \\
\multicolumn{1}{|>{\columncolor[HTML]{C4DCB8}}c|}{0.70} & \multicolumn{1}{|>{\columncolor[HTML]{E5C3B8}}c|}{0.68} & \multicolumn{1}{|>{\columncolor[HTML]{E5C3B8}}c|}{0.68} & \multicolumn{1}{|>{\columncolor[HTML]{D2D2B8}}c|}{0.69} & \multicolumn{1}{|>{\columncolor[HTML]{F5B8B8}}c|}{0.67} & \multicolumn{1}{|>{\columncolor[HTML]{EAC0B8}}c|}{0.68} & \multicolumn{1}{|>{\columncolor[HTML]{D2D2B8}}c|}{0.69} & \multicolumn{1}{|>{\columncolor[HTML]{E5C3B8}}c|}{0.68} \\ \hline
\end{tabular}
%  \end{minipage}
%  \caption{TF exclusion table as described in \autoref{sec:exclusion}. Case $Q_y=0.5$, $D=1$ (left) and $D=2$ (right).}
%  \label{tab:exclusion_Qy_0.5_mw_0_depth_3}
% \end{table}
%  \begin{table}[h]
%  \centering
%  \begin{minipage}{0.1\linewidth}
%   \centering
%   \begin{tabular}{| c|c|}\hline
\textbf{Depth} & \multicolumn{1}{|>{\columncolor[HTML]{B8E6B8}}c|}{Baseline: } \\
 & \multicolumn{1}{|>{\columncolor[HTML]{B8E6B8}}c|}{0.70} \\ \hline
1 & \multicolumn{1}{|>{\columncolor[HTML]{C1DEB8}}c|}{-FOSL1:} \\
 & \multicolumn{1}{|>{\columncolor[HTML]{C1DEB8}}c|}{0.70} \\ \hline
2 & \multicolumn{1}{|>{\columncolor[HTML]{D8CDB8}}c|}{-STAT2:} \\
 & \multicolumn{1}{|>{\columncolor[HTML]{D8CDB8}}c|}{0.69} \\ \hline
3 & \multicolumn{1}{|>{\columncolor[HTML]{F5B8B8}}c|}{-RELA:} \\
 & \multicolumn{1}{|>{\columncolor[HTML]{F5B8B8}}c|}{0.69} \\ \hline
\end{tabular}
%  \end{minipage}
%  \hfill
%  \begin{minipage}{0.85\linewidth}
%   \centering
%   \begin{tabular}{||c|c|c|c|c|c|c|c|}\hline
\multicolumn{8}{|>{\columncolor[HTML]{B8E6B8}}c|}{Baseline: } \\
\multicolumn{8}{|>{\columncolor[HTML]{B8E6B8}}c|}{0.74} \\ \hline
\multicolumn{4}{|>{\columncolor[HTML]{CAD8B8}}c|}{-FOSL1:} & \multicolumn{4}{|>{\columncolor[HTML]{C8D9B8}}c|}{-STAT2:} \\
\multicolumn{4}{|>{\columncolor[HTML]{CAD8B8}}c|}{0.73} & \multicolumn{4}{|>{\columncolor[HTML]{C8D9B8}}c|}{0.73} \\ \hline
\multicolumn{2}{|>{\columncolor[HTML]{D6CFB8}}c|}{-RELA:} & \multicolumn{2}{|>{\columncolor[HTML]{ECBEB8}}c|}{-STAT2:} & \multicolumn{2}{|>{\columncolor[HTML]{ECBEB8}}c|}{-FOSL1:} & \multicolumn{2}{|>{\columncolor[HTML]{CAD8B8}}c|}{-RELA:} \\
\multicolumn{2}{|>{\columncolor[HTML]{D6CFB8}}c|}{0.72} & \multicolumn{2}{|>{\columncolor[HTML]{ECBEB8}}c|}{0.71} & \multicolumn{2}{|>{\columncolor[HTML]{ECBEB8}}c|}{0.71} & \multicolumn{2}{|>{\columncolor[HTML]{CAD8B8}}c|}{0.73} \\ \hline
\multicolumn{1}{|>{\columncolor[HTML]{D6CFB8}}c|}{-FOS:} & \multicolumn{1}{|>{\columncolor[HTML]{F5B8B8}}c|}{-STAT2:} & \multicolumn{1}{|>{\columncolor[HTML]{F5B8B8}}c|}{-RELA:} & \multicolumn{1}{|>{\columncolor[HTML]{ECBEB8}}c|}{-SOX2:} & \multicolumn{1}{|>{\columncolor[HTML]{F5B8B8}}c|}{-RELA:} & \multicolumn{1}{|>{\columncolor[HTML]{ECBEB8}}c|}{-SOX2:} & \multicolumn{1}{|>{\columncolor[HTML]{D6CFB8}}c|}{-FOS:} & \multicolumn{1}{|>{\columncolor[HTML]{F5B8B8}}c|}{-FOSL1:} \\
\multicolumn{1}{|>{\columncolor[HTML]{D6CFB8}}c|}{0.72} & \multicolumn{1}{|>{\columncolor[HTML]{F5B8B8}}c|}{0.70} & \multicolumn{1}{|>{\columncolor[HTML]{F5B8B8}}c|}{0.70} & \multicolumn{1}{|>{\columncolor[HTML]{ECBEB8}}c|}{0.71} & \multicolumn{1}{|>{\columncolor[HTML]{F5B8B8}}c|}{0.70} & \multicolumn{1}{|>{\columncolor[HTML]{ECBEB8}}c|}{0.71} & \multicolumn{1}{|>{\columncolor[HTML]{D6CFB8}}c|}{0.72} & \multicolumn{1}{|>{\columncolor[HTML]{F5B8B8}}c|}{0.70} \\ \hline
\end{tabular}
%  \end{minipage}
%  \caption{TF exclusion table based as described in \autoref{sec:exclusion}. Case $Q_y-0.6$, $D=1$ (left) and $D=2$ (right).}
%  \label{tab:exclusion_Qy_0.6_mw_0_depth_3}
% \end{table}